\SourcePage{170}{175}
\section{Electron in the field of a plane electromagnetic wave}

The Dirac equation can be solved exactly for an electron moving in the field of a plane electromagnetic wave (\textit{D.\,M.\,Volkov}, 1937).

The field of a plane wave with wave 4-vector $k$ ($k^2 = 0$) depends on the 4-coordinates only through the combination $\varphi = kx$, so that the 4-potential
\begin{equation}
A^\mu = A^\mu(\varphi),
\tag{40.1}
\end{equation}
which satisfies the Lorentz gauge condition
\begin{equation}
\partial_\mu A^\mu = k_\mu A'^\mu = 0
\tag{40.2}
\end{equation}
(the prime denotes differentiation with respect to $\varphi$). Since the constant term in $A$ is inessential, in this gauge condition one may drop it and write it in the form
\[
kA = 0.
\]

\SourcePage{171}{176}
We start from the second-order equation~(32.6), in which the field tensor
\begin{equation}
F_{\mu\nu} = k_\mu A'_\nu - k_\nu A'_\mu.
\tag{40.3}
\end{equation}
When expanding the square $(i\partial - eA)^2$ one should take into account that by virtue of~(40.2) $\partial_\mu(A^\mu\psi) = A^\mu\partial_\mu\psi$. As a result we obtain the equation
\begin{equation}
[-\partial^2 - 2ie(A\partial) + e^2A^2 - m^2 - ie(\gamma k)(\gamma A')]\psi = 0
\tag{40.4}
\end{equation}
($\partial^2 = \partial_\mu\partial^\mu$).

We seek the solution of this equation in the form
\begin{equation}
\psi = e^{-ipx}F(\varphi),
\tag{40.5}
\end{equation}
where $p$ is a constant 4-vector. Adding to $p$ any vector of the form $\mathrm{const}\cdot k$ does not change the form of the function $\psi$ (only a corresponding redefinition of the function $F(\varphi)$ is needed). One may therefore impose on $p$ one additional condition without loss of generality. Let
\begin{equation}
p^2 = m^2.
\tag{40.6}
\end{equation}
Then when the field is switched off the quantum numbers $p^\mu$ become the components of the 4-momentum of a free particle. The meaning of the components of the 4-vector $p$ in the presence of the field is more transparent in a special reference frame, chosen so that $A_0 = 0$. Let in this frame the vector $\mathbf{A}$ be directed along the $x^1$ axis, and $\mathbf{k}$ along the $x^3$ axis (i.e.\ the electric field of the wave is polarised along $x^1$, the magnetic along $x^2$, and the wave propagates along $x^3$). Then~(40.5) takes the form of eigenfunctions of the operators
\[
\widehat p_1 = i\frac{\partial}{\partial x^1},\qquad \widehat p_2 = i\frac{\partial}{\partial x^2},\qquad \widehat p_0 - \widehat p_3 = i\!\left(\frac{\partial}{\partial x_0} - \frac{\partial}{\partial x^3}\right)
\]
with eigenvalues $p_1$, $p_2$, $p_0 - p_3$ (the same operators, as is easily verified, are commutative with the Hamiltonian of the Dirac equation). Thus, in this reference frame, $p^1$, $p^2$ are the components of the generalised momentum along the axes $x^1$, $x^2$, and $p^0 - p^3$ is the difference between the total energy and the component of the generalised momentum along the axis $x^3$.

On substituting~(40.5) into~(40.4) we note that
\[
\partial^\mu F = k^\mu F',\qquad \partial_\mu\partial^\mu F = k^2 F'' = 0,
\]
and find for $F(\varphi)$ the equation
\[
2i(kp)F' + [-2e(pA) + e^2A^2 - ie(\gamma k)(\gamma A')]F = 0.
\]
The formal solution of this equation is
\[
F = \exp\left\{-i\int_0^{kx}\!\left[\frac{e}{(kp)}\,(pA) - \frac{e^2}{2(kp)}\,A^2\right]\!d\varphi + \frac{e(\gamma k)(\gamma A)}{2(kp)}\right\}\frac{u}{\sqrt{2p_0}},
\]
where $u/\sqrt{2p_0}$ is an arbitrary constant bispinor (its form will be discussed below).

\SourcePage{172}{177}
All powers of $(\gamma k)(\gamma A)$ above the first vanish, since
\[
(\gamma k)(\gamma A)(\gamma k)(\gamma A) =
\]
\[
= -(\gamma k)(\gamma k)(\gamma A)(\gamma A) + 2(kA)(\gamma k)(\gamma A) = -k^2A^2 = 0.
\]
One can therefore replace
\[
\exp\frac{e(\gamma k)(\gamma A)}{2(kp)} = 1 + \frac{e}{2(kp)}\,(\gamma k)(\gamma A),
\]
so that $\psi$ takes the form
\begin{equation}
\psi_p = \left[1 + \frac{e}{2(kp)}\,(\gamma k)(\gamma A)\right]\frac{u}{\sqrt{2p_0}}\,e^{iS},
\tag{40.7}
\end{equation}
where\,\footnote{The expression for $S$ coincides with the classical action for a particle moving in the field of a wave (cf.~II, \S\,47, problem~2).}
\begin{equation}
S = -px - \int_0^{kx}\!\left[\frac{e}{(kp)}\,(pA) - \frac{e^2}{2}\,A^2\right]\!d\varphi.
\tag{40.8}
\end{equation}

To establish the conditions imposed on the constant bispinor $u$, one should consider the wave as infinitely slowly ``switched on'', starting from $t = -\infty$. Then $A\to 0$ at $kx\to-\infty$ and $\psi$ should go over into the solution of the free Dirac equation. For this $u = u(p)$ must satisfy the equation
\begin{equation}
(\gamma p - m)u = 0.
\tag{40.9}
\end{equation}
This condition discards the ``spurious'' solutions of the second-order equation. Since it does not depend on time, this condition remains in force at finite $kx$ as well. Thus, $u(p)$ coincides with the bispinor amplitude of a free plane wave; we shall take it normalised by the same condition~(23.4): $\bar u\,u = 2m$.

The arguments set forth also immediately allow one to find the normalisation of the wave functions~(40.7). An infinitely slowly switched on field does not change the normalisation integral. It follows that the functions~(40.7) satisfy the same normalisation condition
\begin{equation}
\int\psi^*_{p'}\,\psi_p\,d^3x = \int\bar\psi_{p'}\gamma^0\psi_p\,d^3x = (2\pi)^3\delta(\mathbf{p}' - \mathbf{p}),
\tag{40.10}
\end{equation}
as free plane waves.

Let us find the current density corresponding to the functions~(40.7). Noting that
\[
\bar\psi_p = \frac{\bar u}{\sqrt{2p_0}}\!\left[1 + \frac{e}{2(kp)}\,(\gamma A)(\gamma k)\right]e^{-iS},
\]
by direct multiplication we get
\begin{equation}
j^\mu = \bar\psi_p\gamma^\mu\psi_p = \frac{1}{p_0}\!\left\{p^\mu - eA^\mu + k^\mu\!\left(\frac{e(pA)}{(kp)} - \frac{e^2A^2}{2(kp)}\right)\right\}.
\tag{40.11}
\end{equation}

If $A^\mu(\varphi)$ are periodic functions and their average (over time) vanishes, the mean value of the current density is
\begin{equation}
\bar\jmath^\mu = \frac{1}{p_0}\!\left(p^\mu - \frac{e^2\overline{A^2}}{2(kp)}\,k^\mu\right).
\tag{40.12}
\end{equation}

\SourcePage{173}{178}
Let us find also the mean density of kinetic momentum in the state $\psi_p$. The operator of kinetic momentum is the difference $\widehat p - eA = i\partial - eA$. By direct computation we find
\[
\psi^*_p(\widehat p^\mu - eA^\mu)\psi_p = \bar\psi_p\gamma^0(\widehat p^\mu - eA^\mu)\psi_p =
\]
\begin{equation}
= p^\mu - eA^\mu + k^\mu\!\left(\frac{e(pA)}{(kp)} - \frac{e^2A^2}{2(kp)}\right) + k^\mu\,\frac{ie}{8(kp)p_0}\,F_{\mu\nu}(u^*\sigma^{\lambda\nu}u).
\tag{40.13}
\end{equation}
The time-averaged value of this 4-vector, which we denote by $q^\mu$, is
\begin{equation}
q^\mu = p^\mu - \frac{e^2\overline{A^2}}{2(kp)}\,k^\mu.
\tag{40.14}
\end{equation}
Its square:
\begin{equation}
q^2 = m_*^2,\qquad m_* = m\sqrt{1 - \frac{e^2}{m^2}\,\overline{A^2}}\,;
\tag{40.15}
\end{equation}
$m_*$ plays the role of ``effective mass'' of the electron in the field. Comparing~(40.14) and~(40.12), we see that
\begin{equation}
\bar\jmath^\mu = q^\mu/p_0.
\tag{40.16}
\end{equation}

Let us note also that the normalisation condition~(40.10), expressed with the aid of the vector $\mathbf{q}$, has the form
\begin{equation}
\int\psi^*_{p'}\,\psi_p\,d^3x = (2\pi)^3\frac{q_0}{p_0}\,\delta(\mathbf{q}' - \mathbf{q})
\tag{40.17}
\end{equation}
(the transition from~(40.10) to~(40.17) is most easily carried out in the special reference frame indicated above).
